% !TeX root = ../../Book3.tex
% 纲要标题：### 24.2 模与层
% 纲要位置：计划纲要.md，第 2742 行。

\section{模与层}
\label{b3:sec:2402}

把模的元素局部化到每个素点后，还须要求这些局部元素由分式在邻域内相容地表示。由此得到的伴随层保留张量积与正合列，并把模的局部性质转为层的局部性质。\cref{b3:thm:affine-quasicoherent-equivalence}从全局模恢复仿射空间上的层，本节末的黏合命题则从相容的局部层构造全局对象。层化、逆像、拉回和直像在定理前给出定义；它们在仿射图上的表达也用于下一节的 Proj 构造。


% 纲要内容（仅作写作计划，尚非教材正文）：
% * 模对应的拟凝聚层；仿射模层等价及取全局截面的正合性
% * 局部有限秩自由层、线性过渡矩阵与代数向量丛；截面层及线性态射对应
% * 局部化相容性和黏合

\subsection{模对应的拟凝聚层}
\label{b3:subsec:2402:01}

环 \(A\) 在谱上形成结构层，\(A\) 模 \(M\) 也应在每个点保留 \(M_{\mathfrak p}\)。这里的局部元素不是标量函数，但同样可以在分母不消失的区域内表示为分式。

\begin{definition}\label{b3:def:associated-module-sheaf}
设 \(A\) 为交换环，\(M\) 为 \(A\) 模，\(X=\operatorname{Spec}A\)。对开集 \(U\subseteq X\)，令 \(\widetilde M(U)\) 为所有族 \((s_{\mathfrak p})_{\mathfrak p\in U}\)、\(s_{\mathfrak p}\in M_{\mathfrak p}\) 的集合，要求每个 \(\mathfrak p\in U\) 都有开邻域 \(V\subseteq U\) 以及 \(m\in M,f\in A\)，使所有 \(\mathfrak q\in V\) 都满足 \(f\notin\mathfrak q\) 及 \(s_{\mathfrak q}=m/f\)。配以限制与逐点模运算，所得 \(\mathcal O_X\) 模层称为 \(M\) 的伴随层。这里一个 \(\mathcal O_X\) 模层 \(\mathcal F\) 要求每个 \(\mathcal F(U)\) 为 \(\mathcal O_X(U)\) 模，且限制与标量限制相容。
\end{definition}

\cref{b3:lem:localization-sheaf-gluing}本来就是对模证明的。用同样的局部族构造和分母核验，得到
\begin{equation}\label{b3:eq:module-sheaf-localization}
\Gamma(D(f),\widetilde M)=M_f,\qquad
(\widetilde M)_{\mathfrak p}=M_{\mathfrak p},\qquad
\Gamma(X,\widetilde M)=M.
\end{equation}
模同态在各局部化上给出相容同态，因而诱导层态射。层的核可逐开集取核；余核则一般须把局部余类重新黏合，不能直接把每个开集上的余核当作层。等价地，层序列的正合性在每个茎上检验：一个芽属于像，表示在某个足够小的邻域中有原像，不要求全局已经选好原像。

\begin{definition}\label{b3:def:quasicoherent-sheaf}
设 \(X\) 为概形，\(\mathcal F\) 为 \(\mathcal O_X\) 模层。若存在仿射开覆盖 \(X=\bigcup_iU_i\)、\(U_i=\operatorname{Spec}A_i\)，以及 \(A_i\) 模 \(M_i\)，使 \(\mathcal F|_{U_i}\cong\widetilde{M_i}\)，则称 \(\mathcal F\) 为\term[niningju-ceng]{拟凝聚层}[quasi-coherent sheaf]。在 Noether 概形上，若还能选各 \(M_i\) 为有限生成模，则称 \(\mathcal F\) 为\term[ningju-ceng]{凝聚层}[coherent sheaf]。
\end{definition}

后一用语在本章只用于 Noether 概形，不把“有限生成”作为任意底空间上凝聚性的完整定义。拟凝聚性首先按某个仿射覆盖定义；下一定理说明在仿射空间上，它仍由一个全局模决定。

\ProseParagraphBreak
先说明定理中所用的层运算。对拓扑空间上的预层 \(\mathcal P\)，取所有局部由 \(\mathcal P\) 的同一截面表示的芽族，得到层 \(\mathcal P^+\)，称为其层化。限制和黏合逐芽进行，正如结构层的构造；\(\mathcal P\rightarrow\mathcal P^+\) 在茎上为同构。若 \(\mathcal P\rightarrow\mathcal G\) 的目标已经是层，各局部代表的像便能唯一黏合，故该映射唯一经过 \(\mathcal P^+\)。这样，模层 \(\mathcal F,\mathcal G\) 的张量积就是预层 \(U\mapsto\mathcal F(U)\otimes_{\mathcal O_X(U)}\mathcal G(U)\) 的层化，其茎为 \(\mathcal F_x\otimes_{\mathcal O_{X,x}}\mathcal G_x\)。

\ProseParagraphBreak
对连续映射 \(f:Y\rightarrow X\)，逆像层 \(f^{-1}\mathcal F\) 是预层
\[
U\longmapsto\varinjlim_{f(U)\subseteq V}\mathcal F(V)
\]
的层化，其中 \(U\subseteq Y\)、\(V\subseteq X\) 开，转移映射为缩小 \(V\) 的限制。它的茎为 \((f^{-1}\mathcal F)_y=\mathcal F_{f(y)}\)：两边的代表都是 \(f(y)\) 邻域上的截面，而连续性允许把该邻域拉回到 \(y\) 的邻域。若 \(f\) 是概形态射，就可进一步扩张标量，定义
\(f^*\mathcal F=\mathcal O_Y\otimes_{f^{-1}\mathcal O_X}f^{-1}\mathcal F\)，称为模层的拉回。直像层则直接定义为 \((f_*\mathcal G)(V)=\mathcal G(f^{-1}V)\)；开覆盖的原像仍为开覆盖，所以它已经满足层公理。

\ProseParagraphBreak
一个闭点可以直接区分逆像与拉回。设 \(\mathfrak m\) 为 \(A\) 的极大理想，\(i:\operatorname{Spec}\kappa(\mathfrak m)\to\operatorname{Spec}A\) 为相应态射。单点空间上的层由唯一的茎确定，故
\[
i^{-1}\widetilde M=M_{\mathfrak m},\qquad
i^*\widetilde M
=\kappa(\mathfrak m)\otimes_{A_{\mathfrak m}}M_{\mathfrak m}
\cong M_{\mathfrak m}/\mathfrak mM_{\mathfrak m}.
\]
第一式仍以 \(A_{\mathfrak m}\) 为标量环，保留点附近的芽；第二式通过剩余域扩张标量，只保留该点的纤维值。例如 \(A=k[x]\)、\(\mathfrak m=(x)\)、\(M=A\) 时，前者为 \(k[x]_{(x)}\)，后者为 \(k\)：逆像并没有自动把在点处消失的函数模去。

\begin{theorem}\label{b3:thm:affine-quasicoherent-equivalence}
设 \(A\) 为交换环，\(X=\operatorname{Spec}A\)。函子 \(M\mapsto\widetilde M\) 与 \(\mathcal F\mapsto\Gamma(X,\mathcal F)\) 给出 \(A\) 模和 \(X\) 上拟凝聚层的互逆等价。它们保持正合列；对任意 \(A\) 模 \(M,N\)，有
\[
\widetilde{M\otimes_A N}\cong
\widetilde M\otimes_{\mathcal O_X}\widetilde N.
\]
对任意交换 \(A\) 代数 \(B\)，由 \(A\rightarrow B\) 诱导的 \(f:\operatorname{Spec}B\rightarrow X\) 满足 \(f^*\widetilde M\cong\widetilde{B\otimes_AM}\)。
\end{theorem}
\begin{proof}
先证明拟凝聚层的截面可以清分母。将其定义中的覆盖细化为有限个环境谱 \(X\) 的基本开集 \(D(g_i)\)，使 \(\mathcal F|_{D(g_i)}=\widetilde{N_i}\)。具体地，若 \(D_X(g)\subseteq U=\operatorname{Spec}B\) 且 \(\mathcal F|_U=\widetilde N\)，令 \(b\in B\) 是全局函数 \(g\) 在 \(U\) 上的限制。一个点处的 \(g\) 可逆与其限制 \(b\) 可逆等价，所以 \(D_X(g)=D_U(b)\)，而限制层来自 \(N_b\)。这证明所需细化保持模的描述；有限性来自 \(X\) 的拟紧性。

若 \(s\in\Gamma(X,\mathcal F)\) 限制到 \(D(f)\) 后为零，则它在各 \(N_i\) 中的像被 \(f\) 的某个幂零化。取统一的幂后，由层的唯一性有 \(f^Ns=0\)。若 \(t\in\Gamma(D(f),\mathcal F)\)，则在每个 \(D(g_i)\cap D(f)\) 上，它是 \(N_i\) 的分式。乘统一的 \(f^N\) 可延拓为 \(D(g_i)\) 上截面 \(t_i\)。在每个交集 \(D(g_ig_j)\) 上，差 \(t_i-t_j\) 限制到 \(D(fg_ig_j)\) 为零。对该仿射交集使用刚证的零化结论，再乘统一的 \(f^L\)，各截面即相容，能够黏合为全局截面。

因此，令 \(M=\Gamma(X,\mathcal F)\)，自然映射
\(M_f\rightarrow\Gamma(D(f),\mathcal F)\) 既单又满。它们在基本开集上相容，遂给出 \(\widetilde M\cong\mathcal F\)。式\eqref{b3:eq:module-sheaf-localization}给出另一方向的复合为恒等。层态射取全局截面后，其在基本开集上的作用由分式唯一确定，所以态射也得到互逆对应。

局部化正合，故模的正合列在取伴随层后仍正合。为证明逆函子的正合性，给定拟凝聚层的短正合列
\[
0\longrightarrow\mathcal F'\longrightarrow\mathcal F
\longrightarrow\mathcal F''\longrightarrow0,
\]
令 \(M'=\Gamma(X,\mathcal F')\)、\(M=\Gamma(X,\mathcal F)\)、\(M''=\Gamma(X,\mathcal F'')\)。上面得到的自然同构把它们在每个基本开集 \(D(f)\) 上的截面映射识别为
\(M'_f\rightarrow M_f\rightarrow M''_f\)。再取茎，层列的正合性给出
\[
0\longrightarrow M'_{\mathfrak p}\longrightarrow M_{\mathfrak p}
\longrightarrow M''_{\mathfrak p}\longrightarrow0
\quad\text{对每个 }\mathfrak p\in\operatorname{Spec}A\text{ 成立}.
\]
因此全局模列 \(0\rightarrow M'\rightarrow M\rightarrow M''\rightarrow0\) 在各素点局部化后正合。局部化保持核、像及余核，故这条模列的各个同调模在全部素点为零，因而本身为零：若某个模中有非零元素，取包含其零化子的极大理想，就能在该处局部化后仍检测到它。这证明全局模列正合，也就得到逆函子的正合性。张量积公式在每个基本开集上就是局部化与张量积的相容式。

最后，拉回模层按“先取局部截面的逆像，再把标量由 \(f^{-1}\mathcal O_X\) 扩为 \(\mathcal O_{\operatorname{Spec}B}\)”定义，即
\(f^*\mathcal F=\mathcal O_{\operatorname{Spec}B}\otimes_{f^{-1}\mathcal O_X}f^{-1}\mathcal F\)；在仿射基本开集上代入分式，恰得到 \(B\otimes_AM\) 及其局部化。这证明最后一式。
\end{proof}

例如 \(\widetilde{A/I}\) 在点 \(\mathfrak p\) 的茎为 \(A_{\mathfrak p}/IA_{\mathfrak p}\)，故它的支集是 \(V(I)\)，且仍保留商环中可能存在的幂零元。上述等价的标准叙述及清分母引理见\cite[Chapter II, Lemma 5.3, Proposition 5.4 and Corollary 5.5, pp. 112--113]{Hartshorne1977}。

\subsection{局部自由层与向量丛}
\label{b3:subsec:2402:02}

\begin{definition}\label{b3:def:locally-free-sheaf}
设 \(X\) 为概形，\(\mathcal E\) 为 \(\mathcal O_X\) 模层。若每点都有开邻域 \(U\)，使 \(\mathcal E|_U\cong\mathcal O_U^r\)、\(r\) 为有限非负整数，则称 \(\mathcal E\) 为有限秩\term[jubu-ziyou-ceng]{局部自由层}[locally free sheaf]。秩可在不同连通分量改变；恒为一的情形称为可逆层。
\end{definition}

\begin{corollary}\label{b3:cor:vector-bundle-projective-module}
设 \(A\) 为交换环，\(X=\operatorname{Spec}A\)。在\cref{b3:thm:affine-quasicoherent-equivalence}中，有限生成投射 \(A\) 模恰对应有限秩局部自由层。
\end{corollary}
\begin{proof}
若 \(M\) 有限生成投射，\cref{b3:thm:finite-presentation-flat-projective}使其在每点的某个基本开邻域上为有限秩自由模，式\eqref{b3:eq:module-sheaf-localization}随即给出局部自由层。反之，局部自由层显然拟凝聚，写成 \(\widetilde M\)。取有限个基本开集使其平凡，则每个 \(M_{f_i}\) 为有限秩自由模。\cref{b3:prop:finite-presentation-open-cover}先保证 \(M\) 有限展示，再由同一局部自由判据得 \(M\) 投射。
\end{proof}

纤维与茎保存的信息不同。在点 \(x\) 处，纤维是
\(\mathcal E(x)=\mathcal E_x\otimes_{\mathcal O_{X,x}}\kappa(x)\)；
茎则是邻域数据的局部环模。一个有限生成模的所有纤维维数相同，仍不能在有幂零元的基底上保证局部自由。例如 \(A=k[\varepsilon]/(\varepsilon^2)\)、\(M=A/(\varepsilon)\) 只有一个一维纤维，却不自由：若自由，维数要求它为秩一，而其非零零化子 \((\varepsilon)\) 排除了这一点。

\ProseParagraphBreak
概形本身也可以沿开子概形黏合。给定 \(X_i\)、开集 \(U_{ij}\subseteq X_i\) 和同构 \(\alpha_{ij}:U_{ij}\rightarrow U_{ji}\)，要求 \(U_{ii}=X_i\)、\(\alpha_{ii}=1\)、\(\alpha_{ji}=\alpha_{ij}^{-1}\)，并且
\[
\alpha_{ij}(U_{ij}\cap U_{ik})=U_{ji}\cap U_{jk},\qquad
\alpha_{jk}\circ\alpha_{ij}=\alpha_{ik}
\quad\text{于 }U_{ij}\cap U_{ik}.
\]
先把不交并按这些同构取商；每个 \(X_i\) 到商空间是开嵌入，因为其开集的饱和在各分量仍开。再按相容族的方法黏合环层，茎在每个局部仍是原茎。因此商空间的每点有原来的仿射邻域，得到概形。这个构造同时拼合拓扑和结构层。

\ProseParagraphBreak
局部自由层的几何描述正是这种黏合的一个应用。对开子概形 \(U\) 和非负整数 \(r\)，在各仿射图 \(\operatorname{Spec}A\subseteq U\) 上取
\(\operatorname{Spec}A[T_1,\ldots,T_r]\)。沿 \(U\) 的交叠，用系数的限制及同名变量进行黏合，得到 \(U\) 上的平凡仿射空间，记为 \(\mathbb A_U^r\)；各图上的自然投影也黏合为 \(\mathbb A_U^r\rightarrow U\)。例如在基本开集上，所用限制只是
\(A[T_1,\ldots,T_r]\rightarrow A_f[T_1,\ldots,T_r]\)。共同细化仿射覆盖表明，构造与覆盖选择无关。

\begin{definition}\label{b3:def:algebraic-vector-bundle}
设 \(X\) 为概形，\(p:V\rightarrow X\) 为概形态射。若存在开覆盖 \(X=\bigcup_iU_i\)、有限非负整数 \(r_i\)，以及 \(U_i\) 上的同构
\[
\phi_i:p^{-1}(U_i)\xrightarrow{\sim}\mathbb A_{U_i}^{r_i},
\]
使每个非空交集 \(U_i\cap U_j\) 上有 \(r_i=r_j\)，且坐标变换 \(\phi_i\phi_j^{-1}\) 形如
\(v_i=G_{ij}v_j\)，其中 \(G_{ij}\) 是该交集上正则函数组成的可逆矩阵，则称 \(p\) 及这组相容的线性平凡化为 \(X\) 上的局部有限秩\term[daishu-xiangliangcong]{代数向量丛}[algebraic vector bundle]。两组平凡化之间的过渡仍为可逆线性变换时，视为同一个向量丛结构。向量丛态射是底空间 \(X\) 上的概形态射，要求它在双方同时平凡化的开集上由正则函数矩阵给出，不含常数项；两个丛的秩不同时，矩阵可以是矩形的。
\end{definition}

矩阵条件描述的是基底环上的线性变换，不能只检查每个剩余域纤维上的点映射；在含幂零元的基底上，后者未必检测到坐标表达中的全部项。纤维秩在每个平凡图上为 \(r_i\)，允许不同连通分量有不同秩。

\begin{proposition}\label{b3:prop:locally-free-vector-bundle-equivalence}
设 \(X\) 为概形。局部有限秩代数向量丛与有限秩局部自由 \(\mathcal O_X\) 模层互相对应，且这种对应保持态射的方向。具体地，向量丛 \(p:V\rightarrow X\) 的截面层定义为
\[
\mathcal S_V(U)=\{\,s:U\rightarrow V\text{ 为概形态射}\mid
p\circ s\text{ 是 }U\hookrightarrow X\,\}.
\]
由平凡化中的坐标加法和标量乘法，\(\mathcal S_V\) 成为有限秩局部自由模层；给定局部自由层 \(\mathcal E\)，可以构造截面层同构于 \(\mathcal E\) 的向量丛。两种构造互为逆，精确到与所给局部识别相容的同构。
\end{proposition}
\begin{proof}
在仿射平凡图 \(U_i=\operatorname{Spec}A_i\) 上，\(\mathbb A_{U_i}^{r_i}\) 的截面由保持 \(A_i\) 的环同态
\(A_i[T_1,\ldots,T_{r_i}]\rightarrow A_i\) 给出，恰由各 \(T_j\) 的像决定。因此截面就是 \(r_i\) 个正则函数。进一步限制到开集并黏合，得到 \(\mathcal S_V|_{U_i}\cong\mathcal O_{U_i}^{r_i}\)。线性过渡矩阵保持这些坐标的加法和标量乘法，使它们黏合为模层结构；概形态射的局部黏合也给出截面层的层公理。

反过来，取 \(\mathcal E\) 的仿射平凡覆盖，令 \(e_i\) 为第 \(i\) 个图上的一行基截面。在交集上写
\(e_j=e_iG_{ij}\)，则 \(G_{ij}G_{jk}=G_{ik}\)。以
\(v_i=G_{ij}v_j\) 把 \(\mathbb A_{U_j}^{r_j}\) 与 \(\mathbb A_{U_i}^{r_i}\) 的对应开子概形识别，余圈关系保证能黏合为 \(V\rightarrow X\)。一个相容坐标族 \((v_i)\) 同时给出该向量丛的截面和层 \(\mathcal E\) 的截面 \(e_iv_i\)，从而得到自然同构 \(\mathcal S_V\cong\mathcal E\)。更换基或细化覆盖，只改变这些局部表示，相容的局部同构给出总空间的同构。从原向量丛的截面层出发，取原平凡化对应的基，所还原的过渡矩阵也与原矩阵相同，因而恢复原向量丛。

最后，在同时平凡化的图上，模层态射由矩阵 \(H_i\) 给出；它在交集上的相容条件恰为
\(H_iG_{ij}=G'_{ij}H_j\)。同一矩阵定义从第一向量丛到第二向量丛的坐标映射，并由这个等式黏合为向量丛态射。反向把向量丛态射与截面复合，得到同一个矩阵的模层态射。局部矩阵也保持恒等映射与复合，故两个对应在态射上互为逆。
\end{proof}

这里的总空间在平凡图上的坐标代数是 \(\operatorname{Sym}(\mathcal E^\vee)\)，其中对偶层由
\(\mathcal E^\vee(U)=\operatorname{Hom}_{\mathcal O_U}(\mathcal E|_U,\mathcal O_U)\) 给出，所以其截面层是 \(\mathcal E\)。固定秩向量丛的局部平凡化与截面层构造见 Hartshorne\cite[Chapter II, Exercise 5.18, pp. 128--129]{Hartshorne1977}；该书对 \(\operatorname{Sym}(\mathcal E)\) 的总空间记法给出的截面层为 \(\mathcal E^\vee\)。下面射影直线上的非平凡可逆层将具体说明，局部基未必能够黏合为全局基。

\subsection{局部化相容性与黏合}
\label{b3:subsec:2402:03}

\begin{proposition}\label{b3:prop:sheaf-gluing}
设 \(X\) 为概形，\(X=\bigcup_iU_i\) 为开覆盖。给定各 \(U_i\) 上的 \(\mathcal O_{U_i}\) 模层 \(\mathcal F_i\) 及同构
\(\theta_{ij}:\mathcal F_j|_{U_i\cap U_j}\rightarrow\mathcal F_i|_{U_i\cap U_j}\)。若 \(\theta_{ii}=1\)，且 \(\theta_{ij}\theta_{jk}=\theta_{ik}\) 在每个三重交集上成立，则存在模层 \(\mathcal F\) 及相容同构 \(\mathcal F|_{U_i}\cong\mathcal F_i\)，并在保留这些同构的意义下唯一。若各 \(\mathcal F_i\) 拟凝聚或局部自由，则 \(\mathcal F\) 也分别具有该性质。
\end{proposition}
\begin{proof}
对开集 \(V\subseteq X\)，定义 \(\mathcal F(V)\) 为全部族
\[
s_i\in\mathcal F_i(V\cap U_i),\qquad
s_i=\theta_{ij}(s_j)\ \text{于 }V\cap U_i\cap U_j.
\]
限制和标量作用逐分量给出。局部相容族先在每个 \(U_i\) 内由层公理黏合，所得截面仍满足交集关系，因而 \(\mathcal F\) 是层。在 \(V\subseteq U_i\) 时，第 \(i\) 分量决定其余分量为 \(\theta_{ji}(s_i)\)；余圈关系保证相容，遂得所需同构。任何其他实现都按这些局部同构给出同一个相容族，得到唯一性。最后两项性质本来就是开局部条件。
\end{proof}

取两份 \(\operatorname{Spec}k[t]\)，沿 \(D(t)\) 用恒等映射黏合，得到带双原点的仿射直线。它的全局函数是
\[
\{(a(t),b(t))\in k[t]^2:a=b\text{ 于 }k[t,t^{-1}]\}
\cong k[t].
\]
若此概形仿射，全局函数环将由\cref{b3:thm:affine-scheme-hom}恢复整个概形；但两个原点在全局函数上的取值理想均为 \((t)\)，不可能在同一个素谱中成为两个不同点。故它不是仿射概形。这个例子说明，知道各仿射块和全局函数仍不足以丢掉黏合数据。
